\subsection*{标量矩阵乘法（不使用 DMA）}

第一个版本用 \co@parallel@ 划分输出，用 \co@.at()@ 逐个读取元素，尚未加入显式的数据搬运。

\begin{lstlisting}[language=croqtile,caption={标量矩阵乘法（不使用 DMA）}]
__co__ s32 [128, 256] matmul(s32 [128, 256] lhs, s32 [256, 256] rhs) {
  s32 [lhs.span(0), rhs.span(1)] output;
  parallel {p, q} by [16, 64] {
    foreach {m, n} in [128 / #p, 256 / #q]
      output.at(p#m, q#n) = 0;
    foreach index = {m, n, k} in [128 / #p, 256 / #q, 256]
      output.at(p#m, q#n) += lhs.at(p#m, k) * rhs.at(k, q#n);
  }
  return output;
}
\end{lstlisting}

图~\ref{fig:f05-fig5_scalar_matmul} 画出了输出的分块网格。

\begin{figure}[htbp]\centering
\includegraphics[width=\linewidth,height=0.56\textheight,keepaspectratio]{f05-fig5_scalar_matmul_light.png}
\caption{[128, 256] 输出按 [16, 64] 网格划分，每块形状为 [8, 4]；图中只画出开头的一部分分块。}
\label{fig:f05-fig5_scalar_matmul}
\end{figure}

可以从输出形状开始，逐步核对并行划分、块内范围与全局坐标。

输出声明从 \co@lhs.span(0)@ 取得 128 行，从 \co@rhs.span(1)@ 取得 256 列。收缩维度出现在两个输入中，但不会成为输出维度。

声明 \co@parallel {p, q} by [16, 64]@ 为同一并行层次提供两个坐标轴。1,024 对索引将输出划分为每块八行四列。

循环将 \co@m@、\co@n@ 和 \co@k@ 绑定为 \co@index@ 的分量。前两者遍历输出块内位置，第三者遍历收缩维度。迭代次数直接来自输出划分：

\begin{itemize}[leftmargin=1.6em]
\item \co@128 / #p = 128 / 16 = 8@：每块的行数。
\item \co@256 / #q = 256 / 64 = 4@：每块的列数。
\item \co@256@：完整的收缩维度。
\end{itemize}

组合行索引 \co@p#m@ 的值为 \co@p * 8 + m@。p 遍历 16 个行分块，m 遍历每块八行，合起来覆盖第 0 到 127 行。列表达式 \co@q#n@ 使用相同规则，每块四列。

对每个输出位置，语句体沿 K 累加 \co@lhs.at(p#m, k) * rhs.at(k, q#n)@，所有贡献进入同一个输出元素。清零循环在同一并行实例中先执行，因此每个累加元素都由拥有它的实例清零。
